\id{МРНТИ 65.29.03}{}

\begin{header}
\swa{М.Е.~Сейсеналы, М.П.~Байысбаева, Д.А.~Шаншарова}{ПОЛБА ҰНЫНЫҢ ФИЗИКО-ХИМИЯЛЫҚ САПАЛЫҚ КӨРСЕТКІШТЕРІН ЗЕРТТЕУ}

М.Е.~Сейсеналы\envelope,
М.П.~Байысбаева,
Д.А.~Шаншарова
\end{header}

\begin{affil}
Алматы Технологиялық Университеті, Алматы, Қазақстан

\corrauthor{Корреспондент-автор: makpal-97\_97@mail.ru}
\end{affil}

Бұл мақалада полба ұнының физика-химиялық сапалық көрсеткіштері «Руно»
және «1127-7» сұрыптары негізінде зерттелді. Зерттеудің мақсаты-полба
дәні мен ұнының химиялық құрамы, биологиялық құндылығы және
технологиялық қасиеттерін салыстырмалы түрде бағалап, оларды тағам
өндірісінде тиімді қолдану мүмкіндіктерін анықтау. Зерттеу нәтижелері
бойынша «Руно» сұрыбының дәніндегі крахмал мөлшері 47,93 \% болып,
«1127-7» сұрыбымен (45,03 \%) салыстырғанда жоғары екені анықталды, бұл
оның энергия құндылығының басым екенін көрсетеді. Ақуыз мөлшері екі
сұрыпта да жоғары деңгейде сақталып, тиісінше 17,24 \% және 17,43 \%
құрады. Глютен мөлшері «1127-7» сұрыбында 35 \%, ал «Руно» сұрыбында 31
\% болды, бұл сұрыптардың технологиялық бағытталуының әртүрлі екенін
дәлелдейді. Ұнның тарту ірілігіне байланысты химиялық құрамының өзгеруі
анық байқалды. Ірі тартылған ұнда күлділік пен талшық мөлшері жоғары
болып, «1127-7» сұрыбында күлділік 4,21 \%, талшық 6,84 \% деңгейіне
жетті. Орташа тартылған ұн қоректік және технологиялық қасиеттердің
оңтайлы теңгерімімен сипатталды, ал ұсақ тартылған ұнда глютен мөлшері
жоғарылағанымен, биологиялық құндылық төмендеді. Аминқышқылдық құрамды
талдау ірі тартылған полба ұнының биологиялық құндылығы жоғары екенін
көрсетті. Әсіресе «1127-7» сұрыбында лизин мен триптофан мөлшерінің
жоғары болуы тағамдық артықшылық береді. Витаминдік құрам бойынша «Руно»
сұрыбында В тобы витаминдері, ал «1127-7» сұрыбында Е витамині (12,81
мг/100 г) басым болды. Желімше көрсеткіштері полба ұнының бидай ұнына
қарағанда жоғары екенін көрсетті: желімше мөлшері 37,45-38,24 \%,
созылғыштығы 21-23 \% аралығында болды. Қорытындылай келе, полба ұны
функционалдық және наубайханалық мақсатта қолдануға жарамды, жоғары
тағамдық құндылыққа ие перспективалы шикізат болып табылады.

{\bfseries Түйін сөздер:} полба ұны, физика-химиялық көрсеткіштер,
аминқышқылдық құрам, витаминдер, желімше сапасы.

\begin{header}
ИССЛЕДОВАНИЕ ФИЗИКО-ХИМИЧЕСКИХ ПОКАЗАТЕЛЕЙ КАЧЕСТВА МУКИ ПОЛБЫ

М.Е.~Сейсеналы\envelope,
М.П.~Байысбаева,
Д.А.~Шаншарова
\end{header}

\begin{affil}
Алматинский Технологический Университет, Алматы, Казахстан,

e-mail: makpal-97\_97@mail.ru
\end{affil}

В данной статье исследованы физико-химические показатели качества
полбяной муки на примере сортов «Руно» и «1127-7». Цель исследования
заключалась в сравнительной оценке химического состава, биологической
ценности и технологических свойств зерна и муки полбы, а также в
определении возможностей их эффективного использования в пищевой
промышленности. Результаты исследования показали, что содержание
крахмала в зерне сорта «Руно» составило 47,93 \%, что выше по сравнению
с сортом «1127-7» (45,03 \%), и свидетельствует о его более высокой
энергетической ценности. Содержание белка в обоих сортах находилось на
высоком уровне и составило соответственно 17,24 \% и 17,43 \%.
Количество клейковины у сорта «1127-7» достигало 35\%, тогда как у сорта
«Руно»-31\%, что указывает на различную технологическую направленность
сортов. Установлено, что изменение степени помола существенно влияет на
химический состав муки. В муке крупного помола отмечалось повышенное
содержание зольности и пищевых волокон, при этом у сорта «1127-7»
зольность достигала 4,21 \%, а содержание клетчатки-6,84\%. Мука
среднего помола характеризовалась оптимальным соотношением пищевой и
технологической ценности, тогда как мука мелкого помола, несмотря на
более высокое содержание клейковины, обладала пониженной биологической
ценностью. Анализ аминокислотного состава показал, что мука крупного
помола из полбы отличается высокой биологической ценностью. Особенно у
сорта «1127-7» отмечено повышенное содержание лизина и триптофана, что
повышает его пищевую значимость. По витаминному составу сорт «Руно» был
более богат витаминами группы В, тогда как сорт «1127-7» отличался
высоким содержанием витамина Е (12,81 мг/100 г). Показатели клейковины
свидетельствуют о преимуществах полбяной муки по сравнению с пшеничной:
содержание клейковины составило 37,45-38,24 \%, а ее растяжимость 21-23
\%. В целом полбяная мука является перспективным сырьем с высокой
пищевой ценностью и может эффективно использоваться для производства
функциональных и хлебобулочных изделий.

{\bfseries Ключевые слова:} мука полбы, физико-химические показатели,
аминокислотный состав, витамины, качество клейковины.

\begin{header}
STUDY OF THE PHYSICOCHEMICAL QUALITY INDICATORS OF SPELT FLOUR

M.E.~Seisenaly\envelope,
M.P.~Baysbayeva,
D.A.~Shansharova
\end{header}

\begin{affil}
Almaty Technological University, Almaty, Kazakhstan,

e-mail: makpal-97\_97@mail.ru
\end{affil}

This article investigates the physicochemical quality indicators of
spelt flour using the ``Runo'' and ``1127-7'' varieties as examples. The
aim of the study was to comparatively evaluate the chemical composi\-tion,
biological value, and technological properties of spelt grain and flour,
as well as to determine the possibilities for their effective use in the
food industry. The research results showed that the starch content in
the grain of the ``Runo'' variety was 47.93\%, which is higher than that
of the ``1127-7'' variety (45.03\%), indicating its greater energy
value. The protein content in both varieties was high, amounting to
17.24\% and 17.43\%, respectively. The gluten content of the ``1127-7''
variety reached 35\%, while that of the ``Runo'' variety was 31\%,
demonstrating differences in the technological orientation of the
varieties. It was established that the degree of milling has a
significant effect on the chemical composition of the flour.
Coarse-milled flour was characterized by a higher ash and dietary fiber
content; in particular, the ash content of the ``1127-7'' variety
reached 4.21\%, while fiber content was 6.84\%. Medium-milled flour
showed an optimal balance between nutritional and technological
properties, whereas fine-milled flour, despite its higher gluten
content, was characterized by reduced biological value. Analysis of the
amino acid composition demonstrated that coarse-milled spelt flour
possesses a high biological value. Especially in the ``1127-7'' variety,
increased levels of lysine and tryptophan were observed, which enhances
its nutritional significance. In terms of vitamin composition, the
``Runo'' variety was richer in B-group vitamins, while the ``1127-7''
variety was distinguished by a high content of vitamin E (12.81 mg/100
g). The gluten characteristics indicate the advantages of spelt flour
compared to wheat flour: gluten content ranged from 37.45 to 38.24\%,
and gluten extensibility was 21-23\%. Overall, spelt flour is a
promising raw material with high nutritional value and can be
effectively used in the production of functional and bakery products.

{\bfseries Keywords:} spelt flour, physicochemical indicators, amino acid
composition, vitamins, gluten quality.

\begin{multicols}{2}
{\bfseries Кіріспе.} Полба (\emph{Triticum spelta L.})-ежелгі бидай
түрлерінің бірі болып табылады және қазіргі кезде функционалдық,
экологиялық таза әрі қоректілігі жоғары өнімдер өндіруде маңызды шикізат
көзі ретінде қарастырылуда {[}1{]}. Полба ұны қазіргі жұмсақ бидай
ұнымен салыстырғанда ақуызға, диеталық талшықтарға, минералды
элементтерге және антиоксиданттық қосылыстарға бай келеді. Осы
ерекшеліктері оның тағамдық құндылығын арттырып, диеталық және
профилактикалық бағыттағы өнімдер өндіруге кең мүмкіндік береді {[}2{]}.

Полба ұнының химиялық құрамында калий, фосфор, магний, кальций, темір,
мырыш және марганец элементтері басым болады {[}1{]}. Бұл минералдар зат
алмасу процестерінде, қан түзілу жүйесінде және ферменттік реакцияларда
маңызды рөл атқарады. Полба дәнінің минералдық құрамы сорт ерекшелігіне,
өсу жағдайына және ұнтақтау дәрежесіне байланысты өзгеріп отыратыны
анықталған. Тұтас ұн үлгілерінде ақ ұнға қарағанда микроэлементтер
мөлшері жоғары деңгейде сақталатыны көрсетілген {[}3{]}.

Полба ұны ақуыздың жоғары мөлшерімен және ерекше желімше құрылымымен
сипатталады. Оның глютендік кешені жұмсақ бидайға қарағанда әлсіздеу
болғанымен, биологиялық құндылығы жоғары болып келеді {[}3{]}. Көмірсу
құрамындағы крахмал фракцияларының ерекшелігіне байланысты полба ұнынан
дайындалған өнімдердің гликемиялық индекс көрсеткіші төмен. Сонымен
қатар, тағамдық талшықтардың едәуір мөлшері ас қорыту жүйесінің жұмысын
жақсартып, қан құрамындағы глюкоза деңгейін реттеуге оң әсер етеді {[}2,
11{]}.

Полба ұнының физика-химиялық және реологиялық қасиеттері де ерекше болып
табылады. Оның су сіңіру қабілеті, желімше түзу қасиеті және желімдену
температурасы дәстүрлі бидай ұнынан өзгешелеу келеді. Ұн бөлшектерінің
өлшемі кішірейген сайын су сіңіру қабілеті артып, тұтқырлығы да
жоғарылайтыны анықталған {[}4{]}. Бұл көрсеткіштер қамырдың құрылымына
және дайын өнімнің көлемі мен кеуектілігіне тікелей әсер етеді.

Полба ұнының наубайханалық қасиеттеріндегі кейбір кемшіліктерді жою
мақсатында биопроцессинг тәсілдері кеңінен қолданылуда. Биопроцессинг
барысында микроорганизмдер мен ферменттердің әсерінен ақуыздардың,
полисахаридтердің және фенолдық қосылыстардың құрылымы өзгеріске
ұшырайды {[}5{]}. Мұндай өңдеуден өткен полба ұны қосылған өнімдердің
көлемі ұлғайып, үгінді құрылымы жұмсарып, кеуектілігі жақсаратыны
анықталған. Сонымен қатар антиоксиданттық белсенділік пен фенолдық
қосылыстардың биоқолжетімділігінің артатыны көрсетілген {[}5{]}.

Полба ұнының биологиялық белсенді қосылыстар құрамында фенолдық
қышқылдар, флавоноидтар және токоферолдар едәуір мөлшерде кездеседі.
Аталған қосылыстар организмдегі тотығу-тотықсыздану процестерін реттеп,
жасушаларды бос радикалдардың зиянды әсерінен қорғауға қабілетті.
Зерттеулер нәтижесінде полба ұнындағы антиоксиданттық белсенділік
қазіргі бидай ұнына қарағанда 1,5-2 есе жоғары екендігі анықталған
{[}6,7{]}. Полба дәніндегі тағамдық талшықтардың негізгі бөлігін
ерімейтін целлюлоза, гемицеллюлоза және бета-глюкандар құрайды. Бұл
компоненттер ішек моторикасын жақсартып, ішек микрофлорасының
тепе-теңдігін сақтауға ықпал етеді. Сонымен қатар тағамдық талшықтардың
холестерин деңгейін төмендетуге және семіздік пен қант диабетінің алдын
алуға оң әсер ететіні дәлелденген {[}8,9{]}.

Полба ұнының липидтік фракциясы да жоғары биологиялық құндылыққа ие.
Оның құрамында қанықпаған май қышқылдары, әсіресе линоль және олеин
қышқылдары басым келеді. Бұл май қышқылдары жүрек-қантамыр жүйесінің
қызметін жақсартып, атеросклероздың даму қаупін төмендетеді {[}10{]}.
Полба дәнінің аминқышқылдық құрамы да ерекше назар аударарлық. Лизин,
метионин, треонин және валин сияқты алмастырылмайтын аминқышқылдары
мөлшері жұмсақ бидаймен салыстырғанда жоғары болып келеді. Бұл полба ұны
негізінде дайындалған өнімдердің ақуыздық биологиялық құндылығын
арттырады {[}11{]}.

Полба ұнының технологиялық қасиеттері оны әртүрлі тағам өндірісінде,
атап айтқанда нан, макарон, кондитерлік және диеталық өнімдер
өндірісінде кеңінен қолдануға мүмкіндік береді. Дегенмен таза полба
ұнынан дайындалған қамырдың газ ұстау қабілеті төмендеу болғандықтан, ол
көбінесе жұмсақ бидай ұнымен қоспа түрінде пайдаланылуда {[}12,13{]}.
Соңғы жылдары полба ұнын экструзиялық өңдеу, ультрадыбыстық әсер ету,
ферменттік модификациялау сияқты заманауи технологиялар көмегімен өңдеу
бағыттары қарқынды дамуда. Мұндай өңдеулер полба ұнының реологиялық
қасиеттерін жақсартып қана қоймай, оның биологиялық белсенді заттарының
сақталуын арттырады {[}14,15{]}.

Жалпы алғанда, полба ұны жоғары тағамдық, биологиялық және технологиялық
құндылыққа ие перспективалы шикізат көзі болып табылады. Оның химиялық
құрамы мен физика-химиялық қасиеттерін терең зерттеу отандық
функционалдық және профилактикалық тағам өнімдерін өндіру үшін ғылыми
негіз қалыптастыруға мүмкіндік береді {[}16{]}.

{\bfseries Материалдар мен әдістер.} Бұл зерттеуде негізгі шикізат ретінде
Қазақ егіншілік және өсімдік шаруашылығы ғылыми-зерттеу институтында
селекциялық әдістер арқылы алынған полба (\emph{Triticum spelta L.})
дәндерінің екі түрлі селекциялық үлгісі пайдаланылды. Олар:

- {\bfseries «Руно» сұрыпы}-жоғары бейімделгіштік қасиеттерімен, тұрақты
өнімділігімен және биохимиялық құрамының тұрақтылығымен ерекшеленетін
перспективалы линия;

- {\bfseries «1127-7» сұрыпы}-зертханалық және өндірістік сынақтарда жоғары
нан пісіру сапасын көрсеткен селекциялық материал.

Аталған сұрыптар талдауға дейін стандартталған жағдайларда сақталып,
бастапқы өңдеу сатыларынан (тазалау, қоспалардан арылту, калибрлеу және
уақытша сақтау) өткізілді. Дәндер зертхана жағдайында +18.+20 °C
температурада және салыстырмалы ылғалдылығы 60-65 \% болатын ортада
сақталды.
\end{multicols}

\begin{figs}[1-сурет. \emph{- а - ұсақ тартылған ұн, b - орташа тартылған ұн,
c - ірі тартылған ұн}]
\fig[0.32\textwidth]{p2/image8}[a]
\fig[0.32\textwidth]{p2/image9}[b]
\fig[0.32\textwidth]{p2/image10}[c]
\end{figs}

\begin{multicols}{2}
Зерттеу нысаны ретінде полба ұны алынды. Полба дәндерінен ұн алу процесі
зертханалық жағдайда Chopin Technologies SD-1 Laboratory Mill
диірменінде жүргізілді. Ұнтақтау алдында дәндер қоспалардан тазартылып,
ылғалдылығы 13-14 \% деңгейіне дейін тұрақтандырылды. Диірменнің жұмыс
режимі 10 000 айн/мин ротор жылдамдығында жүзеге асырылды. Ірі дисперсті
фракцияны алу үшін саңылау диаметрі 1000 мкм електер қолданылды; орташа
дисперсті фракция 500 мкм електер арқылы алынды; ал ұсақ дисперсті
фракцияны қалыптастыру үшін 250 мкм електер қолданылды. Осылайша
ұнтақтау параметрлерін (ротор жылдамдығы, електер диаметрі және өңдеу
ұзақтығы) реттеу арқылы бөлшек өлшемі бойынша ірі, орташа және ұсақ
дисперсті полба ұнының фракциялары қалыптастырылды. Алынған ұн бөлшектік
құрамы бойынша шартты түрде үш фракцияға бөлінді:

- ірі тартылған ұнтақ (ірі бөлшектер);

- орташа тартылған ұнтақ;

- ұсақ тартылған ұнтақ (ұнтақталған фракция).

Полба ұнының аминқышқылдық құрамын анықтау үшін капиллярлы электрофорез
әдісі ГОСТ 32195-2013 стандартына сәйкес жүргізіледі. Бұл стандарт тағам
өнімдеріндегі бос және байланысқан аминқышқылдарының мөлшерін жоғары
дәлдікпен анықтауға арналған капиллярлы электрофорез әдісінің жалпы
талаптарын, сынаманы дайындау және талдау жүргізу тәртібін сипаттайды.

Полба ұны құрамындағы В тобы витаминдерінің мөлшері МЕМСТ 31483-2012
стандарты бойынша, ал Е витаминінің мөлшері МЕМСТ 54634-2011 талаптарына
сәйкес анықталды. Аталған стандарттар зертханалық талдауларды
бірізділікпен жүргізуге және нәтижелерді салыстыруға мүмкіндік береді.

Алынған эксперименттік деректер статистикалық әдістермен өңделді. Әрбір
тәжірибе кемінде үш рет қайталанып жүргізілді. Зерттеу нәтижелері
бойынша орташа арифметикалық мән (x̄) және стандарттық ауытқу (SD)
есептелді. Нәтижелердің сенімділігі \emph{p ≤ 0,05} деңгейінде
бағаланды. Деректер орташа мән ± стандартты ауытқу түрінде берілген (M ±
SD), \emph{n = 3}.

{\bfseries Нәтижелер және талқылау.} Полба дәндерінің, ұнының сапалық
көрсеткіштерін анықтау кезінде жүргізілген зерттеулердің жалпы
нәтижелері мен олардың ғылыми тұрғыдағы маңызын қарастыратын боламыз.
Полба ұнының негізгі физико-химиялық көрсеткіштері бағаланды. Сонымен
қатар, зерттеу нәтижелері полба ұнын қолданудың технологиялық
мүмкіндіктерін және оның функционалдық қасиеттерін жалпы сипаттауға
мүмкіндік береді.

Кесте 1 деректері полба дәнінің «Руно» және «1127-7» сорттарының
физика-химиялық сапа көрсеткіштерінде айтарлықтай айырмашылықтар бар
екенін көрсетеді. «Руно» сортының крахмал мөлшері 47,93 \% болып,
«1127-7» сортына (45,03 \%) қарағанда жоғары екені байқалады, бұл оның
энергия құндылығының басым екенін көрсетеді. Ал ылғалдылық көрсеткіші
«1127-7» сортында (9,59 \%) сәл жоғары, бұл оның сақтау тұрақтылығына
белгілі бір дәрежеде әсер етуі мүмкін. Ақуыз мөлшері екі сортта да
жоғары деңгейде болып, тиісінше 17,24 \% және 17,43 \% құрады, бұл
полбаның жоғары тағамдық және биологиялық құндылығын дәлелдейді.
\end{multicols}

\tcap{1-кесте. Полба дәнінің физико-химиялық сапа көрсеткіштері}
\begin{longtblr}[
  label = none,
  entry = none,
]{
  width = \linewidth,
  colspec = {Q[231]Q[115]Q[92]Q[156]Q[115]Q[108]Q[113]},
  cells = {c},
  cells = {font = \small},
  row{1} = {font=\bfseries\small},
  hlines,
  vlines,
}
Полба дәнінің сұрыптары & {\textbf{Крахмал,}\\\textbf{\%}} & {\textbf{Шөгінді,}\\\textbf{\%}} & Ылғалдылық, \% & {\textbf{Ақуыз,}\\\textbf{\%}} & Глютен,\% & Талшық, \%\\
Полба \textbf{«Руно» сорты} & 47,93±0,023 & 78±0,014 & 9,13±0,045 & 17,24±0,052 & 31±0,022 & 29,05±0,12\\
Полба \textbf{«1127-7» сорты} & 45,03±0,025 & 75±0,019 & 9,59±0,039 & 17,43±0,057 & 35±0,025 & 30,45±0,15
\end{longtblr}

\begin{multicols}{2}
Глютен мөлшері бойынша «1127-7» сорты (35 \%) «Руно» сортына (31 \%)
қарағанда басым екені анықталды, бұл оның нан өнімдері үшін
технологиялық тұрғыдан тиімді екенін көрсетеді. Ал талшық мөлшері
«1127-7» сортында 30,45 \%, «Руно» сортында 29,05 \% деңгейінде болып,
екі сорттың да асқорытуға пайдалы функционалдық қасиеттері жоғары екенін
дәлелдейді. Жалпы алғанда, алынған нәтижелер полба дәнінің екі сортының
да жоғары тағамдық, технологиялық және функционалдық қасиеттерге ие
екенін және оларды ашымал өндірісінде тиімді қолдануға болатынын
көрсетеді.

Кесте 2 деректері полба ұнының химиялық құрамы оның тарту ірілігі мен
сорттық ерекшеліктеріне тікелей тәуелді екенін көрсетті. Жалпы алғанда,
екі сұрыпта да ірі тартылған ұн химиялық және биологиялық тұрғыдан ең
бай нұсқа болып табылды, ал ұсақ тартылған ұн тазалану дәрежесі жоғары
болғанымен, қоректік заттар мөлшері төмен нұсқа ретінде сипатталды.

«Руно» сұрыбында ірі тартылған ұн күлділік (2,42 \%), май (3,41 \%),
талшық (3,38 \%) және ақуыз (17,75 \%) мөлшері бойынша ең жоғары
көрсеткіштерді көрсетті. Бұл оның қабықша бөліктерін көбірек қамтитынын
және функционалдық құндылығының жоғары екенін дәлелдейді. Орташа және
ұсақ тартуда аталған көрсеткіштердің төмендеуі байқалып, ең төмен мәндер
ұсақ тартылған ұнда тіркелді. Осыған байланысты «Руно» сұрыбы үшін ірі
тартылған ұн ең тиімді нұсқа болып табылады.
\end{multicols}

\tcap{2-кесте. Әр түрлі іріліктегі полба ұнының химиялық құрамы}
\begin{longtblr}[
  label = none,
  entry = none,
]{
  width = \linewidth,
  colspec = {Q[150]Q[60]Q[100]Q[100]Q[100]Q[100]Q[100]Q[100]},
  cells = {c},
  cells = {font = \small},
  row{1} = {font=\bfseries\small},
  cell{1}{8} = {font=\bfseries\small},
  cell{2}{1} = {r=3}{},
  cell{5}{1} = {r=3}{font=\bfseries\small},
  hlines,
  vlines,
}
Полба дәнінің сұрыптары & Тарту ірілігі & Күлділігі, \% & {\textbf{Май,}\\\textbf{\%}} & Талшық, \% & Ылғалды\-лық, \% & Ақуыз, \% & Глютен, \%\\
Полба \textbf{«Руно» сұрыпы} & ірі & 2,42±0,015 & 3,41±0,032 & 3,38±0,025 & 12,73±0,12 & 17,75±0,21 & 27,1±0,21\\
 & орташа & 1,34±0,011 & 2,96±0,034 & 0,48±0,021 & 12,06±0,09 & 13,63±0,18 & 24,00±0,18\\
 & ұсақ & 0,66±0,013 & 2,55±0,033 & 0,96±0,027 & 12,93±0,11 & 12,60±0,16 & 23,02±0,16\\
Полба «1127-7» сұрыпы & ірі & 4,21±0,018 & 3,32±0,037 & 6,84±0,023 & 14,64±0,15 & 17,91±0,24 & 16,2±0,24\\
 & орташа & 1,96±0,014 & 2,75±0,035 & 1,51±0,026 & 11,76±0,08 & 14,07±0,19 & 13,1±0,19\\
 & ұсақ & 0,57±0,011 & 2,24±0,031 & 0,87±0,022 & 14,42±0,13 & 11,61±0,15 & 21,1±0,15
\end{longtblr}

\begin{multicols}{2}
«1127-7» сұрыбында бұл үрдіс одан да айқын көрінді. Ірі тартылған ұнда
күлділік (4,21 \%) мен талшық мөлшерінің (6,84 \%) айтарлықтай жоғары
болуы бұл сұрыптың қабықшаға бай екенін және диеталық бағыттағы өнімдер
үшін тиімді екенін көрсетеді. Алайда, осы нұсқада глютен мөлшерінің ең
төмен деңгейде (16,2 \%) болуы оның таза наубайханалық мақсатта
қолданылуын шектеуі мүмкін. Керісінше, ұсақ тартылған «1127-7» ұны
глютен мөлшері бойынша (21,1 \%) ең жоғары мән көрсетті, бұл оны
технологиялық тұрғыдан тиімді етеді, бірақ қоректік заттар мөлшері
төмендеу болып келеді.

Тарту ірілігін салыстыра отырып, ірі тартылған ұн-функционалдық және
диеталық құндылығы жоғары нұсқа, орташа тартылған ұн-қоректік пен
технологиялық қасиеттер арасындағы оңтайлы теңгерім, ал ұсақ тартылған
ұн-технологиялық (глютенге бағытталған) тұрғыдан тиімді, бірақ
биологиялық құндылығы төмендеу нұсқа екені анықталды.

Сұрыптар арасында салыстырмалы талдау «1127-7» сұрыбының күлділік пен
талшық мөлшері бойынша басым екенін, ал «Руно» сұрыбының ақуыз және
глютен көрсеткіштерінің тұрақтылығымен ерекшеленетінін көрсетті.
Осылайша, функционалдық тағам өнімдері үшін «1127-7» сұрыбының ірі
тартылған ұны, ал наубайханалық мақсатта «Руно» сұрыбының орташа және
ұсақ тартылған ұны ең тиімді нұсқа болып табылады.
\end{multicols}

\tcap{3-кесте. Полба ұнының тартылу ірілігіне байланысты аминқышқылдық құрамы}
\begin{longtblr}[
  label = none,
  entry = none,
]{
  cells = {c},
  cells = {font = \small},
  row{1} = {font=\bfseries\small},
  cell{1}{2} = {c=3}{},
  cell{1}{5} = {c=3}{},
  hlines,
  vlines,
}
Көрсеткіш & Полба «Руно» сұрыпы &  &  & Полба «1127-7» сұрыпы &  & \\
Тарту ірілігі & ірі & орташа & ұсақ & ірі & орташа & ұсақ\\
Аргинин,\% & 0,94±0,04 & 0,91±0,03 & 0,76±0,03 & 1,15±0,15 & 0,91±0,04 & 0,67±0,02\\
Цистеин,\% & 0,34±0,01 & 0,33±0,01 & 0,32±0,01 & 0,38±0,02 & 0,34±0,01 & 0,32±0,01\\
Гистидин,\% & 0,40±0,02 & 0,38±0,02 & 0,36±0,01 & 0,47±0,02 & 0,39±0,02 & 0,33±0,01\\
Изолейцин,\% & 0,55±0,02 & 0,49±0,02 & 0,49±0,02 & 0,55±0,03 & 0,50±0,02 & 0,46±0,02\\
Лейцин,\% & 1,14±0,05 & 0,99±0,04 & 0,97±0,04 & 1,13±0,05 & 1,01±0,04 & 0,90±0,03\\
Лизин,\% & 0,56±0,02 & 0,52±0,02 & 0,42±0,02 & 0,66±0,03 & 0,48±0,02 & 0,36±0,01\\
Метионин,\% & 0,26±0,01 & 0,26±0,01 & 0,25±0,01 & 0,25±0,01 & 0,35±0,01 & 0,23±0,01\\
Метион+цистеин,\% & 0,59±0,02 & 0,57±0,02 & 0,56±0,02 & 0,64±0,03 & 0,59±0,02 & 0,55±0,02\\
Фенилаланин,\% & 0,72±0,03 & 0,62±0,02 & 0,62±0,02 & 0,71±0,03 & 0,63±0,02 & 0,59±0,02\\
Треонин,\% & 0,47±0,02 & 0,46±0,02 & 0,42±0,02 & 0,56±0,02 & 0,47±0,02 & 0,39±0,01\\
Валин,\% & 0,71±0,03 & 0,66±0,02 & 0,62±0,02 & 0,80±0,03 & 0,67±0,02 & 0,59±0,02\\
Триптофан,\% & 0,23±0,01 & 0,20±0,01 & 0,17±0,01 & 0,29±0,01 & 0,19±0,01 & 0,18±0,01
\end{longtblr}

\tcap{4-кесте. Полба ұнының витаминдік құрамы}
\begin{longtblr}[
  label = none,
  entry = none,
]{
  cells = {c},
  cells = {font = \small},
  row{1} = {font=\bfseries\small},
  hlines,
  vlines,
}
Витамин түрі & «Руно» сорты, мг/100 г & «1127-7» сорты, мг/100 г\\
B₁ (тиамин) & 0,112 ± 0,022 & 0,106 ± 0,021\\
B₂ (рибофлавин) & 0,05 ± 0,021 & 0,03 ± 0,013\\
B₆ (пиридоксин) & 0,11 ± 0,022 & 0,09 ± 0,018\\
B₃ (ниацин) & 2,75 ± 0,55 & 2,63 ± 0,53\\
B₅ (пантотен қышқылы) & Анықталмаған & Анықталмаған\\
B₉ (фолий қышқылы) & 0,015 ± 0,003 & 0,010 ± 0,002\\
E (токоферол) & 2,48 ± 0,01 & 12,81 ± 0,01
\end{longtblr}

\begin{multicols}{2}
Зерттеу нәтижелері полба ұнының аминқышқылдық құрамы оның тарту
ірілігіне тікелей тәуелді екенін көрсетті. Екі сұрыпта да ірі тартылған
ұн алмастырылмайтын аминқышқылдар бойынша ең бай нұсқа болып табылды, ал
ұсақ тартылған ұн аминқышқылдар мөлшері ең төмен нұсқа ретінде
сипатталды.

«Руно» сұрыбында ірі тартылған ұн аргинин, лейцин, валин және триптофан
мөлшері бойынша басым болды. Бұл фракцияда аминқышқылдардың салыстырмалы
түрде жоғары деңгейде сақталуы дәннің қабықша және алейрон қабатының
көбірек сақталуымен байланысты. Орташа тартылған ұнда көрсеткіштердің
біршама төмендеуі байқалғанымен, аминқышқылдық құрамның тепе-теңдігі
сақталды, бұл оны биологиялық құндылық пен технологиялық қолайлылық
арасындағы оңтайлы нұсқа ретінде сипаттауға мүмкіндік береді. Ұсақ
тартылған «Руно» ұны аминқышқылдар мөлшері бойынша ең төмен деңгейде
болып, биологиялық тұрғыдан тиімділігі төмен нұсқа екені анықталды.

«1127-7» сұрыбында аминқышқылдық құрам айқынырақ ерекшеленді. Ірі
тартылған ұн аргинин, лизин, валин және триптофан мөлшері бойынша екі
сұрып арасында да ең жоғары көрсеткіштерге ие болды, бұл осы нұсқаның
биологиялық тұрғыдан ең басым екенін көрсетеді. Әсіресе лизин мен
триптофан мөлшерінің жоғары болуы бұл сұрыптың тағамдық құндылығын
арттырады, себебі бұл аминқышқылдар дәнді дақылдарда шектеуші болып
табылады. Орташа тартуда аминқышқылдар деңгейі айтарлықтай
төмендегенімен, метионин мөлшерінің салыстырмалы түрде артуы осы
фракцияның белгілі бір технологиялық артықшылықтарын сақтайтынын
көрсетеді. Ұсақ тартылған ұн барлық аминқышқылдар бойынша ең төмен
көрсеткіштерге ие болып, тағамдық тұрғыдан ең тиімсіз нұсқа ретінде
бағаланды.

Екі сұрыпты салыстырмалы талдау нәтижесінде «1127-7» сұрыбының ірі
тартылған ұны аминқышқылдық құрамы бойынша ең жоғары биологиялық
құндылыққа ие нұсқа екені анықталды. Ал «Руно» сұрыбының ірі және орташа
тартылған ұны аминқышқылдық құрамының тұрақтылығымен ерекшеленіп,
практикалық қолдану үшін тиімді нұсқа ретінде бағаланды.Жалпы алғанда,
аминқышқылдық құрамы тұрғысынан ірі тартылған полба ұны - ең басым және
тиімді, орташа тартылған ұн -- теңгерімді, ал ұсақ
тартылған ұн - биологиялық құндылығы төмен нұсқа екені ғылыми
тұрғыда дәлелденді.

Кесте 4 деректері полба ұнының «Руно» және «1127-7» сұрыптарының
витаминдік құрамында айқын айырмашылықтар бар екенін көрсетті, бұл
олардың тағамдық және функционалдық құндылығының әртүрлі бағытта
қалыптасатынын дәлелдейді.

В тобы витаминдері бойынша «Руно» сұрыбы жалпы алғанда басым нұсқа болып
табылады. Атап айтқанда, В₁ (тиамин), В₂ (рибофлавин), В₆ (пиридоксин)
және В₉ (фолий қышқылы) мөлшері «Руно» сұрыпында жоғары деңгейде
анықталды. Бұл витаминдер көмірсу, ақуыз және энергия алмасу
процестерінде маңызды рөл атқаратындықтан, «Руно» сұрыбын күнделікті
тағамдық рационды байыту үшін тиімді деп бағалауға болады. Ниацин (В₃)
мөлшері екі сұрыпта да жоғары және шамалас деңгейде болғанымен, «Руно»
сұрыбында сәл басымдық байқалады, бұл оның метаболизмдік белсенділігін
арттырады.

Ал «1127-7» сұрыбы Е витамині (токоферол) бойынша айқын артықшылық
көрсетті. Бұл сұрыпта токоферол мөлшерінің 12,81 мг/100 г дейін жетуі
оның антиоксиданттық қасиетінің өте жоғары екенін көрсетеді. Осыған
байланысты «1127-7» сұрыбы функционалдық, антиоксиданттық және
профилактикалық бағыттағы өнімдер өндіру үшін ең тиімді нұсқа болып
табылады. Е витаминінің жоғары болуы липидтердің тотығуын баяулатып,
өнімнің сақтау тұрақтылығын арттыруға да ықпал етеді.

В₅ (пантотен қышқылы) екі сұрыпта да анықталмаған, бұл оның мөлшері
қолданылған әдістің сезімталдық шегінен төмен екенін көрсетеді және
сұрыптық айырмашылықты айқындауға мүмкіндік бермейді. Жалпы алғанда,
витаминдік құрам тұрғысынан «Руно» сұрыбы - В тобы витаминдеріне бай,
теңгерімді тағамдық нұсқа, ал «1127-7» сұрыбы-Е витамині бойынша басым,
айқын функционалдық бағыттағы нұсқа болып табылады. Осылайша, полба
ұнының әрбір сұрыбы витаминдік құрамы бойынша өзіндік артықшылықтарға ие
және оларды тағам өнімдерінің мақсатты түрлерін өндіруде тиімді
қолдануға болады.
\end{multicols}

\tcap{5-кесте. Ұн үлгілеріндегі желімшенің мөлшері мен реологиялық қасиеттерінің салыстырмалы сипаттамасы}
\begin{longtblr}[
  label = none,
  entry = none,
]{
  cells = {c},
  cells = {font = \small},
  row{1} = {font=\bfseries\small},
  cell{4}{1} = {font=\bfseries\small},
  hlines,
  vlines,
}
Ұн үлгілері & Желімшенің мөлшері, \% & Желімшенің созылғыштығы,\%\\
Бидай ұны & 35,08±0,84 & 19±0,72\\
Полба \textbf{«Руно» сұрыпы} & 38,24±1,02 & 23±0,83\\
Полба «1127-7» сұрыпы & 37,45±0,91 & 21±0,88
\end{longtblr}

\begin{multicols}{2}
Зерттеу нәтижелері (Кесте 5) бидай ұны мен полба ұнының желімше мөлшері
мен созылғыштығы арасында айтарлықтай айырмашылықтар бар екенін
көрсетті. Кесте мәліметтеріне сәйкес, жұмсақ бидай ұнындағы желімше
мөлшері 35,08 \% болса, бұл көрсеткіш полбаның «Руно» сұрыпында 38,24
\%, ал «1127-7» сұрыпында 37,45 \% деңгейінде анықталды. Яғни, полба
ұнының екі сұрыбында да желімше мөлшері дәстүрлі бидай ұнына қарағанда
жоғары екені байқалады.

Желімшенің созылғыштығы көрсеткіші де полба сұрыптарында анағұрлым
жоғары мән көрсетті. Атап айтқанда, бидай ұнында бұл көрсеткіш 19 \%
болса, «Руно» сұрыбында 23 \%, ал «1127-7» сұрыбында 21 \% деңгейінде
болды. Бұл полба ұнының желімшесі серпімдірек әрі созылуға төзімді
екенін көрсетеді. Мұндай қасиет қамырдың құрылымдық тұрақтылығын
арттырып, нан өнімдерінің көлемінің ұлғаюына және кеуектілігінің
жақсаруына оң әсер етуі мүмкін. «Руно» сұрыбы желімше мөлшері мен оның
созылғыштығы бойынша ең жоғары көрсеткіштерге ие болды. Бұл аталған
сұрыптың наубайханалық қасиеттерінің жоғары екенін дәлелдейді. Ал
«1127-7» сұрыбы да бидай ұнымен салыстырғанда жақсы нәтиже
көрсеткенімен, «Руно» сұрыбынан сәл төмен деңгейде болды. Бұл
айырмашылық сұрыптардың генетикалық ерекшеліктері мен ақуыздық құрамының
әртүрлілігімен түсіндіріледі.

Желімшенің мөлшері мен оның созылғыштығын анықтаудың практикалық маңызы
ұнның технологиялық қасиеттерін кешенді бағалауда ерекше рөл атқарады,
себебі бұл көрсеткіштер қамырдың құрылым түзу қабілетін, газ ұстау
қасиетін және дайын өнімнің сапасын (көлемі, кеуектілігі, серпімділігі)
тікелей сипаттайды. Желімше мөлшері жоғары болған жағдайда қамырдың
беріктігі артып, нан және макарон өнімдерінің құрылымы жақсарады, ал
оның созылғыштығы қамырдың өңдеу кезіндегі икемділігін, қалыптау мен
ашыту процестеріндегі тұрақтылығын анықтайды. Сондықтан аталған
көрсеткіштерді анықтау ұнды мақсатты түрде қолдану бағыттарын (нан
пісіру, кондитерлік немесе макарон өндірісі) негіздеуге, сондай-ақ
шикізат сапасын бақылау мен технологиялық процестерді оңтайландыруға
мүмкіндік береді. Осыған байланысты полба ұнын функционалдық және
диеталық бағыттағы нан өнімдерін өндіруде тиімді шикізат ретінде
қолдануға болатынын дәлелдейді.

{\bfseries Қорытынды.}Жүргізілген зерттеу нәтижелері полба (\emph{Triticum
spelta L.}) ұнының «Руно» және «1127-7» сұрыптарының жоғары тағамдық,
биологиялық және технологиялық құндылыққа ие екенін көрсетті. Полба
дәндерінің физика-химиялық көрсеткіштерін талдау барысында екі сұрыпта
да ақуыз мөлшерінің жоғары деңгейде (17,24-17,43 \%) екені анықталып,
олардың биологиялық құндылығын дәлелдеді. «Руно» сұрыбы крахмал мөлшері
бойынша (47,93 \%) басымдық көрсетсе, «1127-7» сұрыбы глютен мөлшері (35
\%) және Е витаминінің жоғары деңгейімен (12,81 мг/100 г) ерекшеленді.

Ұнның тарту ірілігі оның химиялық құрамына айтарлықтай әсер ететіні
анықталды. Ірі тартылған ұнда күлділік, талшық, май, ақуыз және
алмастырылмайтын аминқышқылдардың мөлшері жоғары болып, оның
функционалдық және диеталық құндылығының басым екенін көрсетті. Ұсақ
тартуға қарай бұл көрсеткіштердің 15-30 \% аралығында төмендеуі
байқалды.

Аминқышқылдық құрамды салыстырмалы талдау «1127-7» сұрыбының лизин,
аргинин және валин мөлшері бойынша жоғары екенін, ал «Руно» сұрыбының
лейцин мөлшерімен ерекшеленетінін көрсетті. Екі сұрыпта да
алмастырылмайтын аминқышқылдардың оңтайлы арақатынасы сақталған.

Полба ұнының желімше мөлшері мен созылғыштығы бидай ұнына қарағанда
жоғары болып, оның наубайханалық қасиеттерінің жақсы екенін дәлелдеді.
Жалпы алғанда, алынған нәтижелер полба ұнын жоғары ақуызды,
функционалдық және диеталық бағыттағы тағам өнімдерін өндіру үшін
перспективалы шикізат ретінде қолдануға болатынын ғылыми тұрғыда
негіздейді.
\end{multicols}

\begin{center}
{\bfseries Әдебиеттер}
\end{center}

\begin{refs}
1. Sinkovič L., Tóth V., Rakszegi M., Pipan B. Elemental composition and
nutritional characteristics of spelt flours and wholemeals//Journal of
Elementology.-2023.-Vol.28(1).- P.27-39. DOI
10.5601/jelem.2022.27.4.2358.

2. Deyalage S. T., House J. D., Thandapilly S. J., Malalgoda M.
Nutritional characteristics and physicochemical properties of ancient
wheat species for food applications//Food Bioscience. - 2024.-Vol.
62:105397. DOI 10.1016/j.fbio.2024.105397.

3. Kulathunga J., Reuhs B. L., Zwinger S., Simsek S. Comparative Study
on Kernel Quality and Chemical Composition of Ancient and Modern Wheat
Species: Einkorn, Emmer, Spelt and Hard Red Spring
Wheat//Foods.-2021.-Vol.10(4): 761. DOI 10.3390/foods10040761.

4. Wojciechowicz-Budzisz A.,Skřivan P., Sluková M., Švec I. et al.
Comprehensive Characterization of Micronized Wholemeal Flours:
Investigating Technological Properties across Various
Grains.//Foods.-2024.-Vol.13(1): 39. DOI 10.3390/foods13010039.

5. Mencin M., Markanovič N., Mikulič Petkovšek et al. Bioprocessed
Wholegrain Spelt Flour Improves the Quality and Physicochemical
Characteristics of Wheat Bread.// Molecules .- 2023.
-Vol.28(8):3428.
\href{https://doi.org/10.3390/molecules28083428}{DOI
10.3390/molecules28083428}.

6. Dinu M., Whittaker A. et al. Ancient wheat species and human health:
Biochemical and clinical implications//Nutritional biochemistry.
-2018.-Vol.52.-P.1-9. DOI 10.1016/j.jnutbio.2017.09.001.

7. Horvat D., Šimić G., Drezner G., Lalić A., Ledenčan, et al. Phenolic
Acid Profiles and Antioxidant Activity of Major Cereal
Crops.//Antioxidants.- 2020.-Vol.9(6):527. DOI
10.3390/antiox9060527.

8. Warechowska, M., Anders, A., Warechowski, J.~et al.~The endosperm
microstructure, physical, thermal properties and specific milling energy
of spelt (Triticum aestivum~ssp.~spelta) grain and flour.~//Sci
Rep.-2023.-Vol.13:3629. DOI 10.1038/s41598-023-30285-9.

9. Kandić V. V. Nicolic, M. Simic, et al.
Spelt wheat (Triticum spelta) and common bread wheat compared for
nutritional contents and functional-technological properties//Chilean
journal of agricultural research. -2023.-Vol.83(2). - P.146-158. DOI
10.4067/s0718-58392023000200146.

10. J.Wang E.,Chatzidimitroiou L. Wood, et al. Effect of wheat species
(Triticum aestivum vs T. spelta), farming system (organic vs
conventional) and flour type (wholegrain vs white) on composition of
wheat Flour-Results of a retail survey in the UK and Germany-2.
Antioxidant activity, and phenolic and mineral content//Food
Chemistry:X.-2020.-Vol.6: 100091.DOI 10.1016/j.fochx.2020.100091.

11. J.Wang, M.Barancki, G.Hasanaliyeva et al. Effect of irrigation,
fertiliser type and variety on grain yield and nutritional quality of
spelt wheat (Triticum spelta) grown under semi-arid conditions//Food
Chemistry.-2021.-Vol.358: 129826.DOI 10.1016/j.foodchem.2021.129826.

12. J.Wang, M.Baranski, R.Korkut, H. Kalee, L.Wood, P.Bilsborrow et al.
Performance of Modern and Traditional Spelt Wheat (Triticum spelta)
Varieties in Rain-Fed and Irrigated, Organic and Conventional Production
Systems in a Semi-Arid Environment; Results from Exploratory Field
Experiments in Crete, Greece.//Agronomy.- 2021.-Vol.11(5):890. DOI
10.3390/agronomy11050890.

13. Skrajda-Brdak M. et al. Phenolic nutrient composition and grain
morphology of winter spelt wheat (Triticum aestivum ssp. spelta)
cultivated in Poland //Quality Assurance and Safety of Crops \&
Foods.-2018.-Vol.10(3):1-12. DOI 10.3920/QAS2018.1267.

14. Rakszegi M., Tóth V., Mikó P. The place of spelt wheat among plant
protein sources//Cereal Science.-2023.-Vol.114:103813. DOI
10.1016/j.jcs.2023.103813.

15. D.Rajnincova, Z. Galova, et al. Comparison of nutritional and
technological quality of winter wheat (Triticum aestivum L.) and hybrid
wheat (Triticum aestivum L. x Triticum spelta L.)//central European
agriculture. -2018.-Vol.19(2). -P.437-452. DOI 10.5513/JCEA01/19.2.2146.

16. G. Frakolaki,V.Giannou, et al. Chemical characterization and
breadmaking potential of spelt versus wheat flour//Cereal
science.-2018.Vol.79.-P.50-56. DOI 10.1016/j.jcs.2017.08.023.
\end{refs}

\begin{info}
\hspace{1em}\emph{{\bfseries Авторлар туралы мәліметтер}}

Сейсеналы М.Е.- PhD докторант, Алматы Технологиялық Университеті,
Алматы, Қазақстан, e-mail: makpal-97\_97@mail.ru;

Байысбаева М.П.-техника ғылымдарының кандидаты, профессор, Алматы
Технологиялық Университеті, Алматы, Қазақстан, e-mail:
meruert\_80@mail.ru;

Шаншарова Д.А.-т.ғ.д., профессор, Алматы Технологиялық Университеті,
Алматы, Қазақстан, e-mail:
dinara.shansharova@mail.ru.

\hspace{1em}\emph{{\bfseries Information about the authors}}

Seisenaly M.E.- PhD doctoral student, Almaty Technological
University, Almaty, Kazakhstan, e-mail: makpal-97\_97@mail.ru;

Baysbayeva M.P.- Candidate of Technical Sciences, Professor, Almaty
Technological University, Almaty, Kazakhstan, e-mail:
meruert\_80@mail.ru;

Shansharova D.A.- Doctor of Technical Sciences, Professor, Almaty
Technological University, Almaty, Kazakhstan, e-mail:
dinara.shansharova@mail.ru.
\end{info}
